\section{Introduction}
\label{sec:intro}

\textbf{Autonomous skill acquisition under epistemic uncertainty represents a foundational benchmark for artificial general intelligence.} When human players encounter an unfamiliar interactive environment, they do not read procedural manuals---they visually parse the scene within hundreds of milliseconds, effortlessly abstracting embodiment, spatial boundaries, and functional affordances from core-knowledge visual priors \citep{spelke2007core}. The ARC-AGI-3 benchmark \citep{arcagi3} formalizes this capability by evaluating agents across uninstructed, sequential game environments using Relative Human Action Efficiency (RHAE)---a metric that penalizes superfluous exploratory actions quadratically \citep{chollet2019measure}. Under the strict rules of the official offline track, candidate agents must resolve undisclosed games within a nine-hour wall-clock limit on a single NVIDIA RTX PRO 6000 GPU without external network connectivity, precluding both human intervention and unconstrained exploratory rollouts.

\textbf{However, deploying language and vision-language reasoning models in this interactive regime reveals a severe representational bottleneck.} While modern reasoning backbones demonstrate strong symbolic deduction when provided with clean declarative state descriptions, their execution collapses when exposed to raw $64{\times}64$ integer coordinate matrices. In preliminary experiments (Section~\ref{sec:motivation}), a state-of-the-art frontier model expended over 200--300 exploratory actions navigating an elementary maze serialized as raw tokens, exhausting its action budget and failing the task. In stark contrast, providing the identical model with an aligned visual rendering and geometric breakdown reduced required actions to 6--12, surpassing human baselines. This divergence demonstrates that policy deduction is rarely the primary failure mode; rather, unparsed token grids systematically destroy topological adjacency, forcing autoregressive attention to reconstruct basic spatial relations at prohibitive token and action costs.

\textbf{To resolve this fundamental bottleneck, we conduct extensive empirical ablations and establish two central insights:}
\begin{enumerate}
    \item \textbf{Perceptual Misinduction}---Serializing spatial environments as discrete token matrices causes reasoning models to induce flawed topological models, triggering catastrophic quadratic action penalties during early exploratory steps.
    \item \textbf{The Measured--Guessed Rule}---Geometric invariants, including lattice reachability, spatial enclosures, midpoint tethers, and dihedral symmetries under $D_4 \times \mathbb{Z}^+$, can be computed deterministically in $\mathcal{O}(HW)$ host code without training, whereas high-level affordance roles remain latent under partial observability and must be framed as probabilistic hypotheses paired with unit-action falsification probes.
\end{enumerate}

\textbf{Building on these insights, we introduce \textsc{Sarbloh}, an autonomous open-weight agent architecture engineered around a deterministic perception engine and speculative code execution.} Rather than delegating visual scene decomposition to costly neural inference, our host-side perception engine executes in approximately 35\,ms per frame on a single CPU core, transforming raw environment states into three synchronized modalities: an object-centric symbolic state with temporal bipartite identity tracking, a structured geometric briefing governed by the Measured--Guessed Rule, and an affordance-annotated spatial map. The agent acts exclusively through an embedded IPython execution kernel, synthesizing algorithmic state models and dispatching multi-step trajectories via a speculative execution primitive (\code{await play()}) that monitors environment invariants and halts immediately upon sensory prediction discrepancies.

\textbf{Evaluated under identical offline hardware constraints, \textsc{Sarbloh} achieves state-of-the-art action efficiency across both public and undisclosed environments.} Across the 25 public ARC-AGI-3 game environments, \textsc{Sarbloh} resolves 143 of 183 levels in 8.2 hours, discovering distinct solutions significantly faster than the first Milestone winner on identical Kaggle hardware and model weights. On the official held-out Kaggle evaluation suite, \textsc{Sarbloh} advanced from an initial baseline of 0.27 to a verified score of \textbf{16.53}---a 13.7$\times$ efficiency improvement over the first Milestone winner.

\textbf{The primary contributions of this work are summarized as follows:}
\begin{itemize}
    \item \textbf{We identify and analyze the visual grounding bottleneck in interactive ARC-AGI-3}, formalizing the \textbf{Perceptual Misinduction} failure mode and empirically deriving \textbf{The Measured--Guessed Rule} through controlled representation ablations (§\ref{sec:motivation}).
    \item \textbf{We propose \textsc{Sarbloh}}, an autonomous agent featuring a zero-parameter, $\mathcal{O}(HW)$ deterministic perception engine, a speculative code-as-action planning runtime, and a monotonic append-only KV-cache architecture sustaining a 98.4\% prefix-cache hit rate (§\ref{sec:method}).
    \item \textbf{We evaluate \textsc{Sarbloh} across public and official hidden benchmarks}, demonstrating a 13.7$\times$ score gain over the first Milestone winner, providing hardware-matched comparative telemetry, and detailing the failure modes that govern interactive open-weight reasoning (§\ref{sec:results}).
\end{itemize}
